\section{A. The finite-field gerbe: bands, local classes, and descent}
\label{sec:A}

\subsection{Conventions and the exact hypothesis}
We write $\overline{\Q}\subset\C$ and fix an embedding into
$\overline{\Q}_p$. The Tate object $\Q(1)$ has weight $-2$.
The Frobenius endomorphism $\pi_X$ of a variety over $\F_q$ is induced
by raising its coordinates to their $q$th powers; its pullback on cohomology
has the usual Weil eigenvalues, and $\pi_{\Q(-1)}=q$.
On the geometric-point fibre functor of the Artin category, this is
\emph{geometric} Galois Frobenius. For local Brauer groups we use the
arithmetic normalization
\[
 \operatorname{inv}_K(L/K,\sigma_{\rm ar},a)
       =\frac{\operatorname{ord}_K(a)}{[L:K]}
 \quad(L/K\text{ unramified}).
\]
Here the cyclic algebra has $uxu^{-1}=\sigma_{\rm ar}(x)$ and
$u^{[L:K]}=a$. In particular the endomorphism algebra of a simple
isocrystal of slope $s$ has invariant $-s$. Reversing the reciprocity
normalization reverses all nonarchimedean invariants below; it does not
change the classes represented by the stated cyclic algebras.
For a gerbe we use arrows $\phi_\sigma:{}^\sigma\omega\to\omega$
and cocycle $c_{\sigma,\tau}=\phi_\sigma{}^\sigma\phi_\tau
\phi_{\sigma\tau}^{-1}$. With this convention the class evaluated
at a character is the \emph{opposite} of the corresponding simple
object's endomorphism-algebra class. Indeed the semilinear operators
on its geometric fibre make it a left module over the crossed product
of $c$, and its descended endomorphisms are the centralizer of that
crossed product. The two central-simple-algebra classes sum to zero.

Our conditional category $\mathcal M_q$ is the category of rational
numerical pure motives over $\F_q$, with all Tate twists and the
commutativity constraint corrected by the K\"unneth projectors so that
its realizations have ordinary, rather than super, tensor symmetry.
We assume the full Tate conjecture for every smooth projective variety
and every product over every finite extension of $\F_q$: the order of
the pole of the zeta function in codimension $r$ is the dimension of
numerical cycle classes, with the corresponding cycle-class
surjectivity and semisimplicity at the eigenvalue $1$. This is the
hypothesis used in \citet[\S1, especially Propositions 1.15 and
1.17]{AMilne1994}; it gives full faithfulness after extension of
coefficients for both prime-to-$p$ and crystalline realizations.
Where a polarization is invoked we additionally specify the Hodge
standard positivity for the abelian motives involved. No assertion
about arbitrary Chow motives follows from these hypotheses.

The category is semisimple and Tannakian, and its geometric fibre
functors have a commutative band $G$. Thus $G$ is a group over $\Q$
even when no $\Q$-valued fibre functor exists. We use $H^i(k,T)$ for
continuous protorus cohomology, with the fpqc interpretation of torsors
and banded gerbes; the finite-level derived-limit exact sequence is
retained throughout. It is not our definition that $H^i$ commutes with
inverse limits.

\subsection{A.1. Correcting the Weil-number group and the rigidification}

\begin{theorem}[The band and its canonical decompositions]
\label{A:band}
\status{Proved under the stated Tate hypothesis; the group calculation is unconditional}
Put
\[
 \Gamma=W(p)=\{\lambda\in\overline\Q^{\times}:\
  p^N\lambda\text{ is integral for some }N\geq0,
  \quad |\tau\lambda|=p^{w(\lambda)/2}\ (\tau),\quad
  w(\lambda)\in\Z\}.
\]
Then $G=D(\Gamma)$, where $D$ is the contravariant equivalence from
Galois modules to groups of multiplicative type. Let $\mu_\infty$
be the group of roots of unity, $\Gamma^0=\ker(w)$, and
$M=\Gamma^0/\mu_\infty$. There are canonical identifications
\begin{align*}
 X^*(G^0)&=\Gamma/\mu_\infty=M\oplus\Z[\sqrt p],\\
 G/G^0&=\widehat\Z,\qquad G^0=H^0\times\Gm,
 \qquad H^0=D(M),\\
 X^*(H)&=\Gamma'=\Gamma/\langle\mu_\infty,p\rangle
            =M\oplus\Z/2\Z,\\
 H&=H^0\times\mu_2.
\end{align*}
The maps $\epsilon:G\to\Gm$ and $i:\Gm\to G$ have character
maps $1\mapsto p$ and $\lambda\mapsto w(\lambda)$, respectively,
and $\epsilon i(t)=t^2$. The restriction $i|_{\mu_2}$ splits $H$.
The group $M$ is a $\Q$-vector space, and complex conjugation acts
as $-1$. Consequently $H^0$ is a protorus, an inverse limit of tori
compact at the real place. It is not an algebraic torus of finite type.
\end{theorem}
\begin{proof}
The Tate full-faithfulness assertion identifies endomorphisms of every
tensor construction on a motive, after extension to $\Q_\ell$, with
the centralizer of its semisimple Frobenius. The stabilizer of these
tensor invariants is the Zariski closure of Frobenius. It is
commutative and diagonalizable. Descent therefore makes the band a
group of multiplicative type. Its characters are the Frobenius
eigenvalues, with their actual $\Gal(\overline\Q/\Q)$ action.

All eigenvalues satisfy the displayed conditions. Conversely take
$\lambda\in\Gamma$ and choose $N$ so that $a=p^N\lambda$ is integral
and its weight $m=w(\lambda)+2N$ is positive. Any $m$th root $b$ of
$a$ is integral: it is integral over the ring of algebraic integers,
which is integrally closed. Every complex conjugate of $b$ has
absolute value $\sqrt p$. Thus $b$ is a Weil $p$-integer of weight
one. Honda--Tate supplies an abelian variety over $\F_p$ with $b$
as a Frobenius eigenvalue. The eigenvalue $b^m p^{-N}=\lambda$
then occurs in a tensor power with a Tate twist. The case of weight
zero integral $a$ can instead be treated by Artin motives, since
Kronecker's theorem makes $a$ a root of unity. This proves the whole
character group, not merely a subgroup of it. This also derives the
motive version of Honda--Tate used in
\citet[Proposition 2.6 and Remark 2.7]{AMilne1994}.

The torsion in $\Gamma$ is exactly $\mu_\infty$. The dual
$D(\mu_n)$ is the constant cyclic group $\Z/n\Z$: its geometric
characters have precisely the cyclotomic Galois action. Passing to
the limit gives the constant profinite quotient $\widehat\Z$.
For $\lambda$ of weight $w$, the formula
\[
 [\lambda]\longmapsto
 \bigl([\lambda p^{-w/2}],w\bigr)
\]
is an isomorphism $\Gamma/\mu_\infty\to M\oplus\Z$.
The sign of the square root and every Galois change of that root
are roots of unity, so the formula is canonical on these quotients
and Galois equivariant. Under it $p$ maps to $(0,2)$.
All the asserted exact sequences, the torsion $\Z/2\Z$ in
$\Gamma'$, and the splitting by $i|_{\mu_2}$ follow.

If $u\in\Gamma^0$, every algebraic $n$th root of $u$ is again
in $\Gamma^0$. The quotient by roots of unity is torsion free:
$u^n$ torsion implies $u$ torsion. It is therefore uniquely divisible.
For every embedding, $\overline{u}=u^{-1}$, giving the action $-1$.
A finite Galois-stable lattice in $M$ is the character lattice of a
compact real torus, so these lattices give the asserted inverse limit.
Finally let $a$ be a root of $X^2-X+p$. The imaginary quadratic field
$\Q(a)$ has $p$ split, including $p=2$ (discriminant $-7$).
The element $u=a/\bar a$ has normalized $p$-adic valuations $1,-1$
and is not torsion. Its rational powers give a copy of $\Q$ in $M$;
thus $M$ is not finitely generated, and $H^0$ is not of finite type.
\end{proof}

\begin{remark}[The transcription is not a group definition]
\label{A:monoid}
\status{Refuted literally; corrected by Theorem~\ref{A:band}}
The transcription of 15:5 reads ``entiers alg\'ebriques'' and
``$\nu\in\mathbf Z^+$''. Those elements form a monoid, not the
character group of a rigid category: $p$ is in it and $p^{-1}$ is not.
We have explicitly replaced that text by its group of fractions.
Both positive and negative integral weights occur. Weight zero does
not imply root of unity once denominators at $p$ are allowed.
Also the arrow transcribed as $i|_{\mu_2}:\mu_2\to H^0$ must land
in $H$; its composite with $H\to\mu_2$ is the identity, so it
cannot land in $H^0$.
\end{remark}

\begin{proposition}[Artin versus Artin--Tate rigidification]
\label{A:rigidification}
\status{Refuted as literally formulated in the prompt; corrected statement proved}
Let $F^0$ be the intrinsic fibre functor on
$\operatorname{Niv}^0(\mathcal M)^0$, the Artin motives. The gerbe
of fibre functors extending this $F^0$, with an identification on
that subcategory, has band $G^0$, not $H$. To obtain band $H$ one
must also fix a trivialization of the Tate line. Equivalently one
extends $F^0$ to the Artin--Tate subcategory by $F^0(\Q(1))=\Q$.
Its group is
\[
 G/H=D(\mu_\infty\oplus\Z p)
          =\widehat\Z\times\Gm.
\]
\end{proposition}
\begin{proof}
For any quotient $G\to A$, the automorphisms of a fibre functor
acting identically on the subcategory $\Rep(A)$ form its kernel.
The Artin characters are precisely $\mu_\infty$, giving kernel
$G^0$. Adding the Tate character adds $p\in\Gamma$ and gives $H$.
The two character subgroups intersect trivially, which proves the
last product decomposition. The manuscript itself writes
$F^0(\Q(1))=\Q$ on 15:3 even though $\Q(1)$ is not Artin; that
line explains the missing enlargement, rather than making it true
that Artin rigidification alone has band $H$.
\end{proof}

\begin{proposition}[Geometric base change and the universal Frobenius]
\label{A:P}
\status{Proved under the stated Tate hypothesis; identification of character modules unconditional}
There is a canonical isomorphism
\[
 \Gamma/\mu_\infty\xrightarrow{\sim}
 W(p^\infty):=\varinjlim_{n\mid n'}
   \bigl(W(p^n),\lambda\mapsto\lambda^{n'/n}\bigr).
\]
Thus $G^0=P$, the Weil-number protorus of Langlands--Rapoport and
Milne, and $H^0=P_0$, its weight-zero factor. The base-change
functor $\mathcal M_p\to\mathcal M_{\overline{\F}_p}$ induces
the inclusion $P\hookrightarrow G$ on bands. The inclusion
$\Gamma\hookrightarrow\overline\Q^\times$ is Galois equivariant
and hence defines a canonical $\mathfrak f\in G(\Q)$.
\end{proposition}
\begin{proof}
Send $\lambda\in W(p)$ to its germ at $n=1$. Its germ vanishes
exactly when some positive power is $1$, giving kernel
$\mu_\infty$. Given a representative $a\in W(p^n)$, choose an
$n$th root $b$. All its conjugates have absolute value $p^{w(a)/2}$,
and a sufficiently large power of $p$ makes it integral. Thus
$b\in W(p)$ and its germ is that of $a$. Choices differ by a root
of unity. This proves the assertion and all equivariance claims.
Over $\F_{p^n}$ Frobenius becomes the $n$th power, so the same
map is exactly the character map of geometric base change.
Evaluation $\lambda\mapsto\lambda$ is a Galois-equivariant
homomorphism into $\overline\Q^\times$, which is the definition
of a $\Q$-point of $D(\Gamma)$.
\end{proof}

\begin{proposition}[Centres and crystalline endomorphisms]
\label{A:endomorphisms}
\status{Proved under the stated Tate hypothesis}
For a simple $N\in\mathcal M_q$ the algebra $A_N=\End(N)$ is a
central division algebra over $E_N=\Q[\mathfrak f_N]$, and
\[
 A_N\otimes\Q_p\simeq
 \End_{F\text{-}\mathrm{Isoc}(\F_q)}(T_pN).
\]
The right side means endomorphisms commuting with crystalline
Frobenius, including its semilinear structure. It does not mean all
endomorphisms of the underlying vector space. If $d_N^2=[A_N:E_N]$,
then $\operatorname{rank}(N)=d_N[E_N:\Q]$.
\end{proposition}
\begin{proof}
Simplicity in a semisimple category gives a division algebra.
Full faithfulness at any $\ell\ne p$ identifies its scalar extension
with the centralizer of semisimple Frobenius. The double-centralizer
theorem says that the centre is $\Q_\ell[\mathfrak f_N]$;
faithful scalar extension gives $E_N$ over $\Q$. The crystalline
identity is the $p$-adic full-faithfulness statement of
\citet[Proposition 1.17]{AMilne1994}, applied to $N$ with itself.
A maximal commutative subalgebra of the centralizer has dimension
$\operatorname{rank}(N)$ over the realization field. A maximal
commutative subalgebra of $A_N$ has dimension $d_N[E_N:\Q]$,
which proves the rank formula (ibid., Proposition 2.8).
For abelian varieties the full-faithfulness input is Tate's
endomorphism theorem, so this portion is unconditional. The proof
of Theorem~\ref{A:band} explains precisely the additional tensor
argument needed to pass from abelian varieties to arbitrary pure
motives under the full Tate hypothesis.
\end{proof}

\subsection{A.2. Crystalline realization and the local classes}

\begin{theorem}[The warning about $\Z$ and $\widehat\Z$]
\label{A:crystalline}
\status{Proved; Lead A2 is false in its proposed nonneutral interpretation over $\F_p$}
Over $\F_p$, an isocrystal is a finite-dimensional $\Q_p$-vector
space with an invertible \emph{linear} operator. Forgetting the
operator is an exact faithful tensor functor to $\Q_p$-vector
spaces. Thus $T_p$, after this forgetting, really is a fibre functor
over $\Q_p$. Its Frobenius action generally does not extend
continuously from $\Z$ to $\widehat\Z$. Over
$\overline{\F}_p$ the realization instead takes values in the
$\Q_p$-linear category $\mathcal R_p$ of isocrystals on
$B=W(\overline{\F}_p)[1/p]$. Its fibre-functor gerbe has band
$D(\Q)$ and class $\beta_p$ characterized by
\[
 \operatorname{inv}_p\bigl(s_*(\beta_p)\bigr)=s\pmod\Z,
 \qquad s\in\Q=X^*(D(\Q)).
\]
It is not neutral over $\Q_p$.
\end{theorem}
\begin{proof}
The Witt Frobenius on $W(\F_p)=\Z_p$ is the identity, which proves
the first assertion. The Tate object $\Q(-1)$ has $F=p$.
A continuous representation of the compact group $\widehat\Z$
on a finite-dimensional $\Q_p$-space has compact image. The
image of the subgroup $\Z$ acting by $p^n$ is unbounded. This
already proves the warning on 15:3, including at $p=2$.
More generally an invertible operator extends continuously to
$\widehat\Z$ precisely when the closure of the cyclic subgroup
it generates is compact; equivalently it preserves a lattice, or
all its eigenvalues are $p$-adic units. Necessity follows by taking
the bounded orbit of a lattice and its inverse; sufficiency follows
from the profiniteness of $\mathrm{GL}_r(\Z_p)$.

Over $\overline{\F}_p$, the Dieudonn\'e--Manin simple isocrystal
of slope $a/b$, $(a,b)=1$, has basis $e_0,\ldots,e_{b-1}$ and
$F(e_i)=e_{i+1}$, $F(e_{b-1})=p^a e_0$, semilinearly for the
arithmetic Witt Frobenius. Its endomorphism algebra is the cyclic
division algebra generated by the unramified degree-$b$ extension
and a cyclic shift $u$ satisfying
\[
 u x u^{-1}=\sigma_{\rm ar}^{-1}(x),\qquad u^b=p^a.
\]
For example the unramified scalar $x$ acts on $e_i$ by
$\sigma_{\rm ar}^i(x)$, which verifies the relation directly.
Its invariant is $-a/b$ with our normalization. The slope simples
and their tensor products give the geometric characters $\Q$;
the opposite-algebra computation in our conventions gives the
displayed gerbe invariant $a/b$. This is
the concrete form of \citet[Theorem 2.9, Proposition 2.14, and
Example 3.28]{AMilne1994}. Slope $1/2$ has nonsplit quaternion
endomorphisms, so no $\Q_p$-fibre functor can exist. This concerns
the geometric-base-field category, and does not contradict the
first paragraph for $\F_p$.
\end{proof}

\begin{proposition}[The Artin torsor and the connecting map]
\label{A:eta}
\status{Proved under the stated Tate hypothesis; the Artin calculation is unconditional}
For the Artin--Tate rigidification of Proposition~\ref{A:rigidification},
let $\eta_p=\operatorname{Isom}^{\otimes}(T_p|_{\mathrm{Artin}},F^0_p)$,
with the right action on the source; extend it by the trivial Tate
factor. Under
\[
 H^1(\Q_p,G/H)=H^1(\Q_p,\widehat\Z)
       =\Hom_{\rm cts}(G_{\Q_p},\widehat\Z),
\]
it is the unramified quotient. Arithmetic local Frobenius maps to
$+1$ when the canonical motivic Frobenius is the generator $+1$.
The opposite isomorphism torsor has value $-1$ and gives
$\xi_p=-\partial(\eta_p^{\mathrm{op}})$ instead.
The class of the rigidified gerbe is
\[
 \xi_p=\partial(\eta_p),\qquad
 \partial:H^1(\Q_p,G/H)\longrightarrow H^2(\Q_p,H).
\]
\end{proposition}
\begin{proof}
For $\Spec\F_{p^n}$ the crystalline algebra is
$W(\F_{p^n})[1/p]$, the unramified degree-$n$ extension of
$\Q_p$; the standard Artin algebra is $\Q_p^n$. The torsor
of tensor identifications is therefore the torsor of its embeddings.
These torsors are compatible as $n$ varies. More precisely, over
$\overline\Q_p$ write $U(a)=(\tau(a))_\tau$ for the splitting
of the unramified algebra. If $F$ is its Witt Frobenius, then
$\sigma_{\rm ar}U=UF$. Thus the right-torsor cocycle
$U^{-1}\sigma_{\rm ar}U$ is $F$, the crystalline realization of
$\mathfrak f$, and has value $+1$. This computes the quotient
with its orientation, not just up to sign. The $\Gm$ factor contributes no
$H^1$ by Hilbert 90.

The full $G$-gerbe has the crystalline fibre functor as a
$\Q_p$-object. Imposing an identification with $F^0$ is the
lifting gerbe of this $G/H$-torsor for
$1\to H\to G\to G/H\to1$. The class of a lifting gerbe is
its connecting class: choose local lifts of the torsor's transition
functions; an arrow $\phi_\sigma:{}^\sigma(T_p,U)\to(T_p,U)$
satisfies $U\phi_\sigma={}^\sigma U$. Its failure to satisfy
the cocycle identity is the connecting $H$-valued $2$-cocycle.
This gives the formula with no replacement
of crystalline cohomology by a continuous Galois representation.
\end{proof}

\begin{proposition}[The real obstruction, with the Serre argument repaired]
\label{A:real}
\status{Proved for abelian motives, and hence under the stated hypothesis for all motives}
The full category over $\F_p$ has no $\R$-valued fibre functor.
At infinity the rigidified gerbe has class
\[
 \xi_\infty=(0,1/2)\in
 H^2(\R,H^0)\times H^2(\R,\mu_2)=0\times\Z/2\Z.
\]
The nonzero class is represented by the real Weil extension,
whose lift $j$ of complex conjugation satisfies $j^2=-1$.
\end{proposition}
\begin{proof}
First, Honda--Tate applied to the Weil number $\sqrt p$ gives a
simple abelian surface over $\F_p$. Its endomorphism division
algebra has centre $\Q(\sqrt p)$, local invariant $1/2$ at
each real embedding, and zero at finite places. At $p$ the formula
is $(1/2)[\Q_p(\sqrt p):\Q_p]=1\equiv0$, so this works also
at $p=2$. The algebra has degree two over its centre and the
motive has rank four. Over $\R$ its algebra is
$\mathbb H\times\mathbb H$. A faithful module for that algebra
has real dimension at least eight, whereas a fibre functor must
give dimension four. This is a contradiction.

This repairs the preliminary reading's use of a supersingular
elliptic curve over $\F_p$. Its \emph{geometric} quaternion
endomorphisms need not be defined over $\F_p$; one cannot simply
put that quaternion algebra into its endomorphism ring there.
The original elliptic proof becomes valid after extension to
$\F_{p^2}$, where Honda--Tate supplies a curve with scalar
Frobenius $p$ and all geometric endomorphisms. An $\R$-fibre
functor over $\F_p$ would extend to this category: its restriction
to all Artin motives is standard since
$\Hom(\Gal(\C/\R),\widehat\Z)=0$, and the resulting split
Artin algebra $\R^2$ supplies either extension-of-constants
idempotent. Its fibre on this elliptic curve would be two-dimensional
and admit the definite quaternion algebra, again impossible.

The odd-line argument on 15:11 is also valid after specifying
polarizability. A neutral real category with this band would have
the real character $\sqrt p$, a rank-one object of odd weight.
Its polarization must be alternating and nondegenerate, impossible
on a one-dimensional real space. Equivalently the Tate character
on the connected band is twice a real character: the disconnected
kernel $H=H^0\times\mu_2$ records exactly this forbidden odd
line. Disconnectedness \emph{by itself} is not an obstruction for
an arbitrary polarizable Artin category; the connected band, the
Tate character, and odd-weight polarization are essential to this
first argument.

The computation $H^2(\R,H^0)=0$ is proved in
Theorem~\ref{A:hasse} below. Since any real fibre functor has the
standard Artin restriction and any one-dimensional real Tate line
can be trivialized, nonneutrality of the full category implies
nonneutrality of the rigidified gerbe. Finally
$H^2(\R,\mu_2)\simeq\Br(\R)[2]=\Z/2\Z$, giving the stated
class. The cocycle $j^2=-1$ is precisely its image under
$\mu_2\hookrightarrow\Gm$; on weight $w$ the realization has
semilinear square $(-1)^w$.
\end{proof}

\subsection{A.3. The protorus Hasse calculation}

\begin{lemma}[Compact real tori satisfy the degree-two Hasse principle]
\label{A:compacthasse}
\status{Proved unconditionally}
If $T/\Q$ is a torus compact over $\R$, with character lattice
$L$, then $\Sha^2(\Q,T)=0$. In fact
\[
 H^2(\Q,T)\xrightarrow{\sim}\bigoplus_v H^2(\Q_v,T).
\]
\end{lemma}
\begin{proof}
Complex conjugation $c$ acts by $-1$ on $L$. Choose a finite
Galois splitting field $E$. Continuity and torsion-freeness give
$H^1(\Q,L)=H^1(\Gal(E/\Q),L)$, since a continuous homomorphism
from the profinite group $G_E$ to the discrete lattice has finite,
and hence zero, image. In the finite image acting on $L$, $c=-1$
is central and $L^{\langle c\rangle}=0$. Inflation--restriction
therefore injects $H^1(\Q,L)$ into $H^1(\R,L)$.
In particular $\Sha^1(\Q,L)=0$.

Apply \citet[Theorem I.4.20(a)]{AMilneADT}, with the set of all
places and $M=L$, so that the Cartier dual is $T$. Its perfect
pairing of $\Sha^1(\Q,L)$ with $\Sha^2(\Q,T)$ proves
injectivity. The last portion of the exact sequence in
Theorem I.4.20(b) has cokernel dual to $L^{G_\Q}$.
This group is zero because already $L^c=0$. Hence the localization
map is surjective as well. Every application is to a finite-rank
lattice; no inverse-limit assertion is hidden in this step.
\end{proof}

\begin{theorem}[The global class really is locally determined]
\label{A:hasse}
\status{Proved unconditionally for the band in Theorem~\ref{A:band}}
For $H^0=D(M)$ one has
\[
 \Sha^2(\Q,H^0)=0,\qquad H^2(\R,H^0)=0.
\]
For a finite place $v$, there is a canonical identification
\[
 H^2(\Q_v,H^0)=\Hom\bigl(M^{G_{\Q_v}},\Q/\Z\bigr).
\]
More precisely, for any cofinal countable increasing sequence of
Galois-stable lattices $L_n\subset M$, and $T_n=D(L_n)$,
\[
 H^2(\Q,H^0)=\varprojlim_n H^2(\Q,T_n)
            =\varprojlim_n\bigoplus_v H^2(\Q_v,T_n).
\]
These equalities are conclusions of the proof, not definitions of
cohomology. The assertion in A.3 that $\Sha^2$ of an arbitrary
protorus is simply a limit of finite-level $\Sha^2$ is not a
valid general rule.
\end{theorem}
\begin{proof}
The field $\overline\Q$ is countable, so such a countable exhaustion
exists. The inclusions of lattices dualize to surjective torus
transition maps. Fix $n$ and a splitting field $E_n$ for $T_n$.
Restriction to $E_n$ kills $H^1(\Q,T_n)$ by Hilbert 90; restriction
followed by corestriction is multiplication by $d=[E_n:\Q]$.
Thus $d$ kills that group. Since $M$ is uniquely divisible, it
contains $(1/d)L_n$. For large $m$ the inclusion $L_n\subset L_m$
factors through $(1/d)L_n$, and the resulting torus map factors
through multiplication by $d$ on a copy of $T_n$. Consequently
$H^1(\Q,T_m)\to H^1(\Q,T_n)$ is zero. The inverse system of
these $H^1$ groups is pro-zero, so its inverse limit and its first
derived inverse limit vanish.

The exact sequence of \citet[Proposition 2.8]{AMilne2003} is
\[
 0\longrightarrow\varprojlim{}^1 H^1(\Q,T_n)
 \longrightarrow H^2(\Q,H^0)
 \longrightarrow\varprojlim H^2(\Q,T_n)\longrightarrow0.
\]
The left term has now actually been killed. A class locally zero
maps to a locally zero class at every finite level and is zero
there by Lemma~\ref{A:compacthasse}; hence it is zero globally.
The displayed full inverse-limit description follows from the same
lemma.

The identical argument works over $\Q_v$, choosing a finite
splitting extension and its degree. Local Tate duality
\citep[Theorem I.2.4]{AMilneADT} identifies
$H^2(\Q_v,T_n)$ with $\Hom(L_n^{G_{\Q_v}},\Q/\Z)$.
Invariants commute with this increasing union because each invariant
element already belongs to a finite-rank stable lattice. Applying
$\Hom(-,\Q/\Z)$ proves the finite-place formula. At infinity
$T_n$ is a product of norm-one circles, whose degree-two group is
zero: in cyclic cohomology the invariant circle quotient by the
norm is trivial, or use the norm sequence
$1\to U(1)\to\operatorname{Res}_{\C/\R}\Gm\to\Gm\to1$
and $H^1(\R,\Gm)=0$. The real $H^1$ inverse system is killed
by the same division argument, giving $H^2(\R,H^0)=0$.
\end{proof}

\begin{theorem}[Explicit global and local classes]
\label{A:class}
\status{Proved unconditionally for the constructed gerbe; motivic identification under the stated hypotheses}
Write $v_p(p)=1$, and for $m\in M$ choose a Weil-unit representative
$u$. The number $v_p(m):=v_p(u)$ is independent of that choice.
There is exactly one class
\[
 \xi_0=(a_p,b_{p,\infty})\in H^2(\Q,H^0)\times H^2(\Q,\mu_2)
\]
with local components
\begin{center}
\begin{tabular}{lll}
\toprule
Place & $H^0$ component & $\mu_2$ component\\
\midrule
$\ell\ne p,\infty$ & $0$ & $0$\\
$p$ & $m\mapsto v_p(m)\pmod\Z$, $m\in M^{G_{\Q_p}}$ & $1/2$\\
$\infty$ & $0$ & $1/2$\\
\bottomrule
\end{tabular}
\end{center}
Here $b_{p,\infty}$ is the quaternion Brauer class ramified
exactly at $p$ and infinity, viewed in $H^2(\Q,\mu_2)$ by Kummer.
The compact component $a_p$ has infinite order. The finite local
components determine both factors of $\xi_0$.
\end{theorem}
\begin{proof}
At each finite level $T_n$, restrict the displayed $p$-adic
functional to $L_n^{G_{\Q_p}}$ and put zero at every other place.
Lemma~\ref{A:compacthasse} gives a unique global class with these
local values. Uniqueness makes the finite-level classes compatible.
The derived-limit calculation in Theorem~\ref{A:hasse} lifts them
to exactly one $a_p$. This supplies existence without a motivic
conjecture. The Brauer reciprocity exact sequence supplies exactly
one $2$-torsion class with invariants $1/2$ at $p,\infty$ and zero
elsewhere; Kummer identifies it with $H^2(\Q,\mu_2)$ since
$\Pic(\Q)=0$. If all finite invariants are prescribed, their
sum fixes the infinite invariant, proving the last assertion.

For comparison with the motivic gerbe, geometric base change gives
the isocrystal slope morphism
\[
 D(\Q)\longrightarrow P_{\Q_p},\qquad
 X^*(P)\longrightarrow\Q,\quad [\lambda]\longmapsto v_p(\lambda).
\]
On the splitting $X^*(P)=M\oplus\Z[\sqrt p]$, its character map
is $(m,n)\mapsto v_p(m)+n/2$. Theorem~\ref{A:crystalline}
therefore gives the compact local functional $v_p$ and the
weight-line invariant $1/2$. At $\ell\ne p$, geometric
\'etale cohomology is the fibre functor extending the standard
Artin--Tate functor, giving neutrality. The real computation is
Proposition~\ref{A:real}. The resulting global class is the unique
one just constructed.

One may check the connecting-map normalization directly. For
$m\in M^{G_{\Q_p}}$, some $u^n=a$ is fixed by local Galois:
$\sigma u/u$ ranges in a finite group of roots of unity. Lifting
the unramified $\widehat\Z$ cocycle through $D(\Gamma^0)$
produces the cyclic class with invariant
$\eta_p(\sigma_{\rm ar})\operatorname{ord}_p(a)/n=v_p(u)$.
This is the same class as $\partial\eta_p$ in
Proposition~\ref{A:eta}, not a second unrelated local prescription.

For infinite order, use $u=a/\bar a$ from the last paragraph of
Theorem~\ref{A:band} and choose the $p$-adic embedding with
$v_p(u)=1$. The class $[u]/n$ lies in $M^{G_{\Q_p}}$ and has
valuation $1/n$. The functional thus takes the values $1/n$
for all $n\geq1$; no nonzero integer annihilates it.
\end{proof}

\begin{proposition}[An explicit arithmetic description of the character module]
\label{A:valuations}
\status{Proved unconditionally}
For finite Galois CM fields $K\subset\overline\Q$, let $X_K$
be their sets of primes above $p$. Then
\[
 M\simeq\varinjlim_K
 \{f:X_K\to\Q:\ f(cw)=-f(w)\},
\]
where transition maps pull back functions along maps of primes,
and $[u]$ maps to $w\mapsto\operatorname{ord}_w(u)/
\operatorname{ord}_w(p)$. In particular the local groups and the
functional in Theorem~\ref{A:class} are specified by explicit
rational permutation modules.
\end{proposition}
\begin{proof}
A weight-zero Weil number generates a CM field unless it is real,
in which case it is $\pm1$. On its Galois closure, complex
conjugation is central because $c(u)=u^{-1}$ holds for every
conjugate. Its valuations are supported at $p$ and are
anti-invariant under $c$. If all are zero, it is an algebraic unit
all of whose conjugates have absolute value one, hence a root of
unity. This proves injectivity.

Conversely scale a rational anti-invariant function by an integer
so that the corresponding exponents, including ramification
indices, are integral. They define a fractional ideal $I$ supported
at $p$ with $cI=I^{-1}$. A positive power $I^h=(b)$ is principal.
Then $b/\bar b$ has exactly twice these scaled valuations and
all its complex absolute values are one. Taking the requisite
root, and then its class modulo roots of unity, recovers the given
function. This operation is compatible with extension of fields
and proves surjectivity.
\end{proof}

\subsection{A.4. The comparison with the pseudomotivic gerbe}

\begin{theorem}[The correct comparison, and the strict-band counterexample]
\label{A:comparison}
\status{Strict equivalence with band $H$ or $G$ refuted; corrected equivalence proved under the stated hypotheses}
Let $\mathscr H$ be the Artin--Tate rigidified gerbe of class
$\xi_0$. Its pushout along $H\hookrightarrow P=G^0$ is the
pseudomotivic $P$-gerbe of
\citet[Theorem 3.29 and Theorem 3.31]{AMilne1994}, with the
inverse-limit and adelic corrections of
\citet[Theorem 6.3 and Definition 6.6]{AMilne2003}.
Equivalently it is the gerbe obtained by rigidifying only the
Artin fibre functor $F^0$, and under the motivic hypotheses it
is the gerbe of motives over $\overline{\F}_p$.

There is no band-preserving equivalence of the $H$-gerbe itself
or of the full $G$-gerbe with that $P$-gerbe. Their bands have
component groups $\mu_2$, $\widehat\Z$, and $1$, respectively.
The morphisms that do exist are pushouts
\[
 \mathscr H\longrightarrow\mathscr P
       \longrightarrow\mathscr G,
 \qquad H\hookrightarrow P\hookrightarrow G.
\]
\end{theorem}
\begin{proof}
The character isomorphism of Proposition~\ref{A:P} identifies the
connected band with $P$, not merely with an abstract isogenous
group. Under $H=H^0\times\mu_2\hookrightarrow
P=H^0\times\Gm$, the local compact class is unchanged and the
$\mu_2$ class becomes the real/quaternion weight class.
At $p$ this is exactly the pushout of $\mathcal R_p^\vee$ by
its slope morphism; at infinity it is exactly the pushout of the
real realization gerbe by the weight cocharacter. At the other
places it has the standard \(\ell\)-adic object. These are the
local morphisms in Milne's definition, with the above signs.
The uniqueness theorem for this \emph{adelically supplied} system
is \citet[Theorem 6.3]{AMilne2003}; its proof uses
Propositions 3.5, 3.6, 3.10, and 4.3, rather than incorrectly
discarding all derived limits. Hence this is an equivalence of
systems with local data, not merely a match of names.

The base-change tensor functor has band inclusion $P\to G$ by
Proposition~\ref{A:P}. Artin descent identifies its image gerbe
with the rigidification of the Artin functor: on any finite
extension the Artin algebra becomes split and its idempotents
recover the extension-of-constants category; compatible
idempotents give the algebraic closure. Adding a trivialization
of the Tate line takes the further fibre of $P\to\Gm$, which
is $\mathscr H$. This proves both identifications of gerbes.

Any equivalence of banded gerbes induces an isomorphism of their
bands. An isomorphism preserves the identity component and its
component quotient. The three quotients just computed are
nonisomorphic. This is a complete counterexample to the literal
claim in Lead A1 and A.4, for every prime $p$, including $2$.
It also proves why the corrected pushout statement is the
appropriate replacement.
\end{proof}

\begin{proposition}[Dictionary between gerbes and groupoids]
\label{A:dictionary}
\status{Proved}
Let $\mathscr P$ be a $P$-gerbe over $\Q$ and choose
$x\in\mathscr P(\overline\Q)$. Its groupoid over
$\Spec(\overline\Q)\times_{\Spec\Q}\Spec(\overline\Q)$
is the sheaf of isomorphisms between the two pullbacks of $x$.
Its diagonal kernel is $P_{\overline\Q}$. For each place $v$,
a realization functor induces a morphism of such groupoids;
its map on kernels is $1\to P_{\Q_\ell}$ for $\ell\ne p$,
$D(\Q)\to P_{\Q_p}$ at $p$, and $\Gm\to P_\R$ at infinity.
Conversely fpqc descent for this groupoid recovers $\mathscr P$.
Changing $x$ changes the groupoid by an isomorphism whose
conjugacy ambiguity is exactly the change of the chosen
identification of objects.
\end{proposition}
\begin{proof}
Composition of isomorphisms supplies multiplication and the
identity of $x$ supplies the unit; the cocycle condition is
associativity of composition. Restriction to the diagonal gives
$\Aut(x)=P_{\overline\Q}$. An object on an arbitrary base is
locally an object isomorphic to a pullback of $x$, together with
its descent isomorphisms. Those isomorphisms are exactly a torsor
for the groupoid satisfying its cocycle identity. This constructs
inverse functors and proves the converse, rather than identifying
a groupoid with its kernel alone. A tensor realization gives a
contravariant morphism on fibre-functor gerbes. The kernel maps
are read from weight and slope, as calculated above, so all
morphisms in the dictionary type-check. This is the dictionary
used in \citet[\S3, Examples 3.26--3.28]{AMilne1994} and
\citet[\S6]{AMilne2003}.
\end{proof}

\subsection{A.5. Unconditional existence and the remaining enrichment issues}

\begin{theorem}[Unconditional $H$-gerbe and its automorphisms]
\label{A:unconditional}
\status{Proved unconditionally as a statement about banded gerbes}
The class in Theorem~\ref{A:class} defines an $H$-gerbe over
$\Q$ with exactly the prescribed local classes. It is unique up
to $H$-equivalence. No nonzero class can be added without changing
some local class. The group of isomorphism classes of its
band-preserving self-equivalences is
\[
 H^1(\Q,H)=\Q^\times/\Q^{\times2}\times
  \frac{\prod_n T_n(\Q)}
  {\{(a_n f_n(a_{n+1})^{-1})_n:(a_n)_n\in\prod_nT_n(\Q)\}},
\]
for the explicit lattice exhaustion in Theorem~\ref{A:hasse};
$f_n:T_{n+1}\to T_n$ is its transition. The group of automorphisms
of any one such equivalence is $H(\Q)=\{\pm1\}$.
\end{theorem}
\begin{proof}
Banded commutative gerbes are classified by $H^2$; the descent
cocycle construction of Proposition~\ref{A:eta} realizes any such
class. Existence and uniqueness were proved in
Theorem~\ref{A:class}. Twisting a band-preserving equivalence by
an $H$-torsor gives every other one, with $H(\Q)$ as its
$2$-automorphism group. This is the elementary torsor form of
the gerbe classification in \citet[\S5]{AMilne2003}.

The same pro-zero argument used for $H^1(\Q,T_n)$ gives
$H^1(\Q,H^0)=\varprojlim{}^1 T_n(\Q)$. The displayed quotient
is the explicit countable-system definition of that derived
limit, and Kummer gives the other factor. Finally any rational
point of $H^0$ evaluates on $m/n$ in a fixed number field
(the stabilizer of $m/n$ is the stabilizer of $m$). Its value on
$m$ is therefore divisible by every positive integer inside that
number field's multiplicative group. The intersection of all
power subgroups there is $1$, by valuations, Dirichlet units,
and the finite root-of-unity group. Thus $H^0(\Q)=1$.
\end{proof}

\begin{proposition}[The bare category and the category with local data]
\label{A:enrichment}
\status{Bare classification proved; classification of arbitrary framed local data by the bare class refuted}
Fix the band diagram $H\hookrightarrow G$, the grading $i$,
and the augmentation $\varepsilon$, with the Artin--Tate
rigidification specified in Proposition~\ref{A:rigidification}.
Then $\xi_0$ determines the underlying Tannakian category with
these structures up to compatible equivalence, and conversely.
It does not in general determine a family of chosen local
neutralizations extending the fixed rigidification.

Write $S=\{\ell:\ell\ne p,\ \ell\text{ finite}\}$. The precise
replacement for that stronger classification is relative
cohomology, defined by the shifted cone of
$R\Gamma(\Q,H)\to\prod_{\ell\in S}R\Gamma(\Q_\ell,H)$.
For a fixed locally trivializable class its fibre is a torsor under
\[
 \coker\left[H^1(\Q,H)\longrightarrow
              \prod_{\ell\in S}H^1(\Q_\ell,H)\right].
\]
The same statement with a restricted adelic product applies when
an integral compatibility condition is imposed.
\end{proposition}
\begin{proof}
The gerbe of fibre functors and the category of its finite-rank
representations recover each other by tensor descent. The band
diagram recovers the specified subcategory and its rigidification;
$i$ recovers the grading and $\varepsilon$ the Tate augmentation.
Equivalence classes of $H$-gerbes are exactly $H^2(\Q,H)$, so
these constructions prove both directions of the bare assertion.
They do not choose objects in any local gerbe.

Once an object of a neutral local gerbe has been chosen, every
other object is obtained by twisting by an $H$-torsor, uniquely
up to its isomorphism class in $H^1(\Q_\ell,H)$. A global
band-preserving equivalence changes this family by the
localizations of a global $H$-torsor. This proves the displayed
cokernel description directly. Equivalently a relative
$2$-cocycle is a global $2$-cocycle together with a local
$1$-cochain whose coboundary is its localization; quotienting
by global and local coboundaries gives exactly that description.

Here is a fully explicit obstruction already at $p=2$. Twist a
fixed neutralization at $3$ by $3\in\Q_3^\times/\Q_3^{\times2}$,
in the $\mu_2$ factor, and leave every other odd-prime
neutralization unchanged. This changes no local $H^2$ class.
If a global equivalence identified the two families, its
$\mu_2$-torsor would be represented by
$a\in\{3,-3,6,-6\}$ modulo rational squares: its odd
valuations are prescribed, and only the valuation at $2$ is
unrestricted. The classes $3$ and $-3$ are nonsquares at $5$;
$6$ is a nonsquare at $7$; and $-6$ is a nonsquare at $13$.
None has the prescribed trivial localizations. The compact
factor cannot repair this, since $H=H^0\times\mu_2$ and the
projection to the second factor is functorial. These residue
checks are certified in the accompanying program. The change
has finite support, so the counterexample also survives the
usual almost-everywhere integral requirement.

This refutes the stronger reading in A.3(c), rather than the
carefully separated assertion actually transcribed on 15:7.
That page first says ``la classe de $\mathcal H$ comme
$H$-gerbe'' for $\mathcal M,\mathrm{grad},\mathrm{aug},F_0$,
then introduces $H^2(S\bmod T,H)$ after including the $F_\ell$.
If the chosen identifications with $F^0$ are forgotten, one
must further quotient by their local automorphism actions; this
is a different moduli problem, and the framed counterexample
does not assert otherwise. For the $P$-gerbe with the coherent
adelic object and the two local realization morphisms, the
specific uniqueness theorem is \citet[Theorem 6.3]{AMilne2003},
with $\Q^\times$ as the residual $2$-automorphism group.
That theorem includes more than the calculation $\Sha^2=0$.
\end{proof}

\begin{proposition}[Forgetting the local comparison framings]
\label{A:unframed-adelic}
\status{Proved for coherent adelic systems with the prescribed Artin restrictions}
Fix the category and global Artin--Tate rigidification of
Proposition~\ref{A:enrichment}. Suppose that the prime-to-$p$
realizations form one object over the finite adeles away from $p$,
and that this adelic object's Artin restriction is isomorphic to
$F^0\otimes\mathbb A_f^p$, not just separately at each place. At $p$
fix the Artin restriction of the ordinary crystalline $\Q_p$
fibre functor, namely the unramified Artin torsor computed in
Proposition~\ref{A:eta}. If the local comparison isomorphisms with
these restrictions are \emph{not} part of the data, all such
coherent systems on the fixed gerbe are equivalent. One may also
include the prescribed real realization morphism.

Thus the underlying class suffices for the equivalence class of
this unframed coherent system. This assertion and the framed
counterexample in Proposition~\ref{A:enrichment} concern different
specified data. The equivalence itself is not asserted to be
unique as an equivalence preserving the global $H$-rigidification.
\end{proposition}
\begin{proof}
Put $A=G/P=\widehat\Z$. Over any one of the finite coefficient
fields, the isomorphism torsor between two full $G$-fibre functors
has class in $H^1(k,G)$. Their Artin restrictions are isomorphic,
so its image in $H^1(k,A)$ is zero. Exactness for
$1\to P\to G\to A\to1$ therefore lifts that class to
$H^1(k,P)$. A chosen comparison of the Artin restrictions effects
this lift; the result after forgetting the comparison does not
depend on that choice. The argument also works over the
prime-to-$p$ adelic coefficient ring: the stipulated coherent
Artin comparison supplies a reduction of the adelic isomorphism
torsor to $P$.

For fixed-band real realization morphisms into the $G$-gerbe,
their difference is a $G$-torsor by the same descent argument for
the sheaf of isomorphisms of morphisms. It lifts to a $P$-torsor
because $H^1(\R,\widehat\Z)=0$: a continuous homomorphism
from the group of order two to the torsion-free group
$\widehat\Z$ is trivial. If no real morphism is specified, choose
the zero local difference there. The resulting local differences
give a class in $H^1(\mathbb A,P)$, with exactly the adelic
meaning used by \citet[\S6.2]{AMilne2003}.

Propositions 3.6 and 3.10 of that source prove
\[
 H^1(\Q,P)\xrightarrow{\sim}H^1(\mathbb A,P).
\]
This is the specific global-to-adelic theorem, not a consequence
of the bare degree-two Hasse principle. It provides a global
$P$-torsor whose localizations give the differences. Moreover
\[
 H^1(\Q,H)\longrightarrow H^1(\Q,P)
\]
is surjective: in $H=H^0\times\mu_2\to P=H^0\times\Gm$
the first factor is the identity and $H^1(\Q,\Gm)=0$.
Lift the global torsor along this map and twist the global
$H$-equivalence by it. Its underlying $G$-equivalence then
identifies every unframed local realization and still preserves
the global rigidification. This proves existence in the required
direction, including the nontrivial Artin torsor at $p$.

There remains at least the ambiguity from
$H^1(\Q,\mu_2)=\Q^\times/\Q^{\times2}$: its image in
$H^1(\Q,P)$ is zero, while it can change the global
$H$-equivalence. Hence uniqueness of its isomorphism class is not
part of this assertion. For the $P$-system itself, the sharper
uniqueness and residual $2$-automorphisms are exactly those of
Milne's Theorem 6.3, already stated above.
\end{proof}

\begin{proposition}[What the Lefschetz category actually supplies]
\label{A:lefschetz}
\status{Proved unconditionally; equality of its band with $P$ is refuted}
Let $\operatorname{LMot}(\overline{\F}_p)$ be the category generated
by abelian varieties, with correspondences generated by divisor
classes and with the ordinary Tannakian commutativity constraint.
It is a semisimple Tannakian category. Its band is
\[
 L=\mathop{\prod}_{\Pi}^{\Gm} L_\Pi,\qquad
 X^*(L_\Pi)=\Z^{\Pi}/
 \{f:c f=f,\ \textstyle\sum_{\pi\in\Pi}f(\pi)=0\},
\]
where $\Pi$ runs through the geometric simple isogeny classes,
identified with Galois orbits of integral weight-one Weil germs.
The fibre product identifies every multiplier
$[\pi]+[c\pi]$. In particular $L$ is commutative, but it is
not the Weil-number protorus $P$ in general. Compatible powers of Frobenius generate
a subgroup $P\hookrightarrow L$. Prime-to-$p$ realization and
crystalline realization, the weight and Tate morphisms, and the
real obstruction are available unconditionally. The claimed
identification of the entire character group with Weil numbers,
full Frobenius-centralizer faithfulness for arbitrary tensor
objects, and the resulting formula for $H^0$ cannot simply be
transferred to this smaller-correspondence category.
\end{proposition}
\begin{proof}
The Tannakian assertion is \citet[Theorem 1.7]{AMilne1999};
Proposition 1.8 identifies the band on the category generated by
an abelian variety $A$ with its Lefschetz group, and Corollary 1.9
passes to the full category. Over the algebraic closure of a finite
field all abelian varieties have many endomorphisms. The displayed
character module is exactly Lemma 4.2, and its fibre product is
Theorem 4.3 there. These are unconditional applications of
Honda--Tate and the divisor theorem, not the full Tate conjecture.
The natural map on characters sends a formal generator indexed by
$\pi$ to its Weil germ; its kernel consists of the additional
relations among actual Frobenius numbers. This defines $P\to L$.
Every Weil germ is generated by weight-one integral germs and
Tate twists, by the root argument in Theorem~\ref{A:band}; hence
the character map is surjective and the group map is injective.
The following example proves that its kernel on characters need
not vanish.
\end{proof}

\begin{example}[An explicit extra Frobenius relation]
\label{A:lefschetz-counterexample}
\status{Refuted: $L=P$ for all finite-field Lefschetz motives}
Take $p=5$, $\alpha=1+2i$, and
$\beta=(3+\sqrt{-11})/2$. They are Frobenius numbers of the
elliptic curves
\[
 E_\alpha:y^2=x^3+x,\qquad
 E_\beta:y^2=x^3+4x+2
 \quad\text{over }\F_5,
\]
which have four and three rational points, respectively. Choose
$\lambda^2=\alpha\beta$. This is an integral weight-one
$5$-Weil number, with irreducible polynomial
\[
 x^8-6x^6+15x^4-150x^2+625.
\]
Honda--Tate supplies a simple abelian fourfold $B/\F_5$ with this
Frobenius number. In
\[
 h^1(B)^{\otimes2}\otimes h^1(E_\alpha)^\vee
                      \otimes h^1(E_\beta)^\vee
\]
there is a geometric eigenline on which Frobenius is $1$, but
the Lefschetz group acts nontrivially. Thus its character is zero
on $P$ and nonzero on $L$.
\end{example}
\begin{proof}
Both quadratic numbers have norm $5$; taking a square root of
their product preserves the required absolute value and
integrality. Elimination gives the displayed polynomial. Exact
factorization over $\Q$ and both finite-field point counts are
recorded by \texttt{check\_lefschetz\_counterexample.py}.
There is also a conceptual degree check: in
$K=\Q(i,\sqrt{-11})$ the valuations of $\alpha\beta$ at the
four split primes over $5$ are $(2,1,1,0)$, so it is not a square.
Let $L_1=K(\lambda)$, of degree eight over $\Q$. At the two
$K$-primes where $v_5(\alpha\beta)=1$, the extension $L_1/K$
is ramified quadratic; the local degree times slope is
$2(1/2)=1$. At the other primes the slopes are $0$ or $1$.
All Honda--Tate invariants over $\F_5$ therefore vanish, giving
division index one and dimension $[L_1:\Q]/2=4$.
The four conjugates of $\alpha\beta$ have distinct valuation vectors; consequently
no quotient of two of them is a root of unity. All powers
therefore still generate $K$. Over $\F_{25}$ its slopes are
$1,1/2,1/2,0$. Honda--Tate gives division index $2$ and dimension
$2[K:\Q]/2=4$. This is still dimension four, so no multiplicity appears on
base change to $\F_{25}$. Over every further even extension the
same slopes and centre persist, so $B$ remains simple geometrically.
Its geometric centre is $K$, distinct from the two elliptic
centres, and its Lefschetz group is the norm-similitude torus
of $K/K^+$. These conclusions also verify every hypothesis of the
fibre-product description in Proposition~\ref{A:lefschetz}.

Over an algebraic closure of the coefficient field, choose the
eigencharacters labelled by $\lambda,\alpha,\beta$. The tensor
character $2[\lambda]-[\alpha]-[\beta]$ evaluates on Frobenius
as $\lambda^2/(\alpha\beta)=1$. For the geometric band we
use Frobenius squared over $\F_{25}$, where it centralizes all
geometric endomorphisms; its value on the character is still $1$.
In the $B$-factor of $L$,
vary one embedding by $t$ and its conjugate by $t^{-1}$,
leave the other conjugate pair fixed, and leave the two elliptic
factors fixed. The common multiplier is $1$, so this is a
cocharacter of $L$. It evaluates the tensor character as $t^2$.
Thus the character is nonzero in $X^*(L)$, as asserted.
\end{proof}

\begin{proposition}[The unconditional gerbes used in Kisin's paper]
\label{A:kisin}
\status{Known construction; the signed dictionary proved and independently checked}
Kisin's quasi-motivic gerbe has a band $Q$ defined by a universal
condition on real and $p$-adic cocharacters. Its pushout by
$Q\twoheadrightarrow P$ is the pseudomotivic gerbe, with the
same $P$ and the same local realization morphisms as above.
This construction is unconditional. It supplies a Tannakian
category of representations of that gerbe; identifying that
category with all numerical motives is a separate claim with the
hypotheses of A.1. The quasi-motivic band $Q$ itself is not a
replacement name for $G$, $H$, or $P$.
\end{proposition}
\begin{proof}
For a finite Galois field $L/\Q$, Kisin defines
\[
 Q^L=(\operatorname{Res}_{L(\infty)/\Q}\Gm
       \times\operatorname{Res}_{L(p)/\Q}\Gm)/\Gm.
\]
Lemma 3.1.4 gives its universal cocharacter condition; the
transition maps of 3.1.5 include their local-degree factors.
At the other finite places the source gerbe is trivial, at
infinity it is the real fundamental extension, and at $p$ it is
the inverse limit of the local fundamental extensions, of band
$D(\Q)$; these are 3.1.6. Theorem 3.1.9 gives existence and
uniqueness, with the specified almost-everywhere integral
condition of 3.1.8 and compatibility of the Galois band action.

For CM fields, 3.5.1--3.5.5 define the quotient $P$ by Weil
characters, with the finite-level integrality condition on the
local norms. That condition gives finite-rank lattices; it
exhausts the Weil-germ group after the CM field is enlarged.
Indeed a fixed rational valuation denominator is cleared by
enlarging the local degree, and its weight remains unchanged.
Lemma 3.5.3 sends both universal cocharacters to the weight and
slope cocharacters of $P$. The definition in 3.5.5 then takes
the pushout of the gerbe. The signs require an explicit
local-source identification. Kisin's equation (3.5.2), with
$|p|_p=p^{-1}$, pairs his $p$-adic cocharacter with a Weil
character by the negative normalized valuation. In 3.3.3 he
chooses a lift $d_\sigma$ of arithmetic Witt Frobenius whose
$n$th power at the unramified degree-$n$ level is $p^{-1}$.
Thus that source fundamental class has invariant $-1/n$ in
our normalization. Our isocrystal gerbe has invariant $+1/n$.
The inversion $[-1]:D(\Q)\to D(\Q)$ identifies the two source
gerbes: it takes $+1/n$ to $-1/n$. Precomposing Kisin's slope
morphism with this inversion changes $-s$ into $+s$.
Equivalently, pushing his negative fundamental class through
his negative valuation cocharacter gives our positive local
class. The target band $P$ is unchanged. At infinity inversion
does not change the order-two Brauer class. Consequently this
is the groupoid/gerbe operation in
Proposition~\ref{A:dictionary}, with the specified source-band
identification, not a silent identification of opposite signs.
Kisin notes explicitly in 3.1.5, footnote 13, that his auxiliary
map has the inverse convention to Langlands--Rapoport's; the
band and local-class dictionary must include this choice.
All these source statements were opened in
\citet[\S\S3.1 and 3.5]{AKisin2017}.

Kisin's application to points of abelian-type Shimura varieties
is Theorem 0.3, expanded as Theorem 4.6.7, with $p>2$ in this
paper and a conjugation ambiguity in the adelic action. It
does not assert the full Tate conjecture for arbitrary varieties
and it does not justify dropping that hypothesis in A.1.
\end{proof}

\begin{remark}[The literature date changes the abelian scope]
\label{A:2026-abelian}
\status{Known from the opened October 2026 theorem; the implication is stated explicitly}
The dated primary preprint of \citet[Theorem I]{AvanHoftenSempliner2026}
proves the rational Hodge conjecture for CM abelian varieties.
Its \S1.1.2--1.1.3 applies Milne's Hodge-to-Tate implication
to all abelian varieties over finite fields. With these cited
results the full-correspondence category of abelian motives,
and all its tensor powers, has the Weil band $P$
unconditionally. This does not change the Lefschetz category:
the latter was defined by restricting its morphisms to those
generated by divisors, and Example~\ref{A:lefschetz-counterexample}
still proves that its band can be larger. Nor does this theorem
settle the Tate conjecture for arbitrary smooth projective
varieties. The preprint status and exact version are retained in
the bibliography rather than represented as an earlier theorem.
\end{remark}
